% ── Chapter 5: Conclusion and Future Work ─────────────────────
\chapter{บทสรุปและงานในอนาคต (Conclusion and Future Work)}

% ─────────────────────────────────────────────────────────────
\section{บทสรุป (Conclusion)}
\label{sec:conclusion}
% ─────────────────────────────────────────────────────────────

\subsection{สรุปผลการดำเนินโครงงาน (Summary of the Project)}

โครงงานนี้ประสบความสำเร็จในการออกแบบ พัฒนา และประเมินผลเว็บแอปพลิเคชันที่ใช้เทคโนโลยีปัญญาประดิษฐ์ (AI) สำหรับวิเคราะห์ผลการทดสอบความไวต่อยาปฏิชีวนะด้วยวิธี Kirby-Bauer disk diffusion (AST) แบบอัตโนมัติ ระบบที่พัฒนาขึ้นช่วยแก้ปัญหาคอขวดสำคัญในกระบวนการทำงานของห้องปฏิบัติการจุลชีววิทยาคลินิก โดยเฉพาะขั้นตอนการวัดขนาดบริเวณยับยั้งเชื้อ (ZOI) ด้วยไม้บรรทัด ซึ่งเป็นงานที่ใช้เวลานาน มีความผันแปรตามดุลยพินิจของผู้ปฏิบัติงาน และยากต่อการตรวจสอบย้อนหลัง \cite{humphries2018} การเปลี่ยนกระบวนการนี้ให้เป็นอัตโนมัติไม่เพียงแต่ช่วยเพิ่มปริมาณงาน (Throughput) ที่ห้องปฏิบัติการรองรับได้ แต่ยังช่วยยกระดับความสามารถในการทำซ้ำของผลลัพธ์ (Reproducibility) ซึ่งส่งผลดีต่อการเลือกยาปฏิชีวนะเพื่อรักษาผู้ป่วยได้อย่างถูกต้องและแม่นยำ ท่ามกลางวิกฤตการณ์การดื้อยาต้านจุลชีพระดับโลก \cite{who_amr2022}

ระบบได้บูรณาการองค์ประกอบทางเทคนิคที่ซับซ้อน 3 ส่วนหลักเข้าเป็น Pipeline แบบ End-to-end เพียงหนึ่งเดียว:

\begin{enumerate}
  \item \textbf{การตรวจจับวัตถุและการวัดขนาด ZOI (Object Detection and ZOI Measurement):} โดยเลือกใช้โมเดล YOLOv11n เป็นแกนหลัก \cite{wang2024yolov11} ควบคู่กับ Faster R-CNN \cite{ren2015faster} และเปรียบเทียบประสิทธิภาพกับวิธี Hough Circle Transform แบบดั้งเดิมเพื่อเป็นเกณฑ์มาตรฐาน (Baseline)
  \item \textbf{การระบุชื่อยาปฏิชีวนะ (Antibiotic Label Recognition):} พัฒนาขึ้นโดยอาศัยเทคโนโลยี EasyOCR \cite{baek2019craft} ร่วมกับ Pipeline การประมวลผลภาพเบื้องต้นแบบ 2 เส้นทาง (Otsu Thresholding และ CLAHE) เพื่อความแม่นยำในการสกัดรหัสยาจากพื้นผิวแผ่นดิสก์ที่มีขนาดเล็ก
  \item \textbf{การปรับเทียบมาตราส่วนแบบพลวัต (Dynamic Calibration):} โดยใช้เส้นผ่านศูนย์กลางจริงขนาด 6 มม. ของแผ่นยาปฏิชีวนะเป็นวัตถุอ้างอิงภายในการถ่ายภาพแต่ละครั้ง เพื่อแปลงหน่วยพิกเซลเป็นมิลลิเมตรได้อย่างแม่นยำ
\end{enumerate}

ผลลัพธ์จากการประมวลผลทั้งหมดถูกนำเสนอผ่านเว็บแอปพลิเคชันที่พัฒนาด้วย React ซึ่งแสดงค่าเส้นผ่านศูนย์กลาง ZOI พร้อมผลการจำแนกความไวต่อยา (S/I/R) ตามมาตรฐาน CLSI 2025 และ EUCAST 2023 โดยผู้ใช้สามารถตรวจสอบ แก้ไข และบันทึกประวัติการวิเคราะห์เข้าสู่ฐานข้อมูลได้อย่างเป็นระบบ

\subsection{สรุปผลเชิงปริมาณ (Quantitative Summary)}

ผลการทดลองเชิงปริมาณชี้ให้เห็นถึงความสำเร็จที่สำคัญ ดังนี้:

ในด้านความแม่นยำในการวัด ZOI โมเดล **YOLOv11n** แสดงประสิทธิภาพที่เหนือกว่าอย่างชัดเจน โดยมีค่า mAP@0.5 สูงถึง **95.91\%** และค่าความคลาดเคลื่อนสัมบูรณ์เฉลี่ย (MAE) ต่ำเพียง **2.5--2.7 มม.** ซึ่งดีกว่าโมเดล Faster R-CNN (MAE 4.65--4.85 มม.) และวิธี Hough Transform (MAE > 5 มม.) อย่างมีนัยสำคัญ สถาปัตยกรรมแบบ Single-stage ของ YOLOv11n ยังให้ค่าความแม่นยำ (Precision) และค่า Recall ที่เกือบสมบูรณ์ ($\approx$1.0) พร้อมการลู่เข้าสู่จุดสมดุล (Convergence) ที่รวดเร็ว

สำหรับการปรับเทียบภาพ ระบบ **Dynamic Calibration** ประสบความสำเร็จในการลดความคลาดเคลื่อนเชิงระบบ (Bias) โดยการคำนวณอัตราส่วนพิกเซลต่อมิลลิเมตรแบบรายแผ่นยาในทุกภาพถ่าย ทำให้ระบบสามารถรองรับความแปรปรวนของระยะถ่ายภาพและความบิดเบี้ยวของเลนส์กล้องที่แตกต่างกันในแต่ละสถานการณ์ได้จริง

สำหรับระบบ OCR การใช้ **EasyOCR ร่วมกับ Otsu Thresholding** บรรลุความแม่นยำในการระบุชื่อยาได้ถึง **95\%** ซึ่งสูงกว่าการใช้ภาพถ่ายดิบหรือโมเดลอื่นๆ อย่างเห็นได้ชัด ยืนยันว่า Pipeline การเตรียมภาพเบื้องต้นเป็นกุญแจสำคัญที่ช่วยให้ AI สามารถอ่านข้อความขนาดเล็กบนพื้นผิวจำกัดได้อย่างมีประสิทธิภาพ

ในด้านความเร็ว ระบบสามารถประมวลผลแบบครบวงจรเสร็จสิ้นภายในเวลาเพียง **3--5 วินาทีต่อภาพ** ซึ่งครอบคลุมทั้งการตรวจจับ, OCR, การแปลผล และการบันทึกฐานข้อมูล ซึ่งนับเป็นการเพิ่มความเร็วในการทำงานได้มากกว่า 100 เท่าเมื่อเทียบกับการวัดด้วยมือแบบดั้งเดิม

\subsection{สถานะความสำเร็จและการติดตั้งใช้งานจริง (Project Status and Deployment)}

ปัจจุบันโครงงานนี้ **ดำเนินการแล้วเสร็จตามวัตถุประสงค์** และได้เปลี่ยนผ่านจากต้นแบบงานวิจัยสู่เครื่องมือสนับสนุนทางคลินิกที่พร้อมใช้งานจริง Pipeline ทั้งหมดได้รับการบูรณาการเข้าด้วยกันอย่างสมบูรณ์แบบ และระบบได้ถูก **ติดตั้งใช้งานจริงบนเครื่องเซิร์ฟเวอร์ของวิทยาลัยแพทยศาสตร์ศรีสวางควัฒน ราชวิทยาลัยจุฬาภรณ์ (CRA)** โดยเปิดให้บุคลากรเข้าถึงผ่านเครือข่ายภายในที่ปลอดภัยที่ \texttt{https://paniti-jupyter.cra.ac.th/zoneanalyzer/} ระบบรองรับตารางจุดตัดมาตรฐานล่าสุดทั้ง **CLSI 2025** และ **EUCAST 2023** พร้อมจับคู่อัตโนมัติ ซึ่งเป็นการยืนยันความพร้อมของเทคโนโลยีในการนำไปช่วยสนับสนุนงานในห้องปฏิบัติการจุลชีววิทยาจริง

% ─────────────────────────────────────────────────────────────
\section{ข้อจำกัด (Limitations)}
\label{sec:limitations}
% ─────────────────────────────────────────────────────────────

เพื่อให้การแปลผลลัพธ์ในบริบททางคลินิกเป็นไปอย่างถูกต้อง จึงควรพิจารณาข้อจำกัดของระบบในปัจจุบัน ดังนี้:

\subsection{ความคลาดเคลื่อนเชิงระบบที่เหลืออยู่ (Residual Bias)}

แม้ว่าระบบการปรับเทียบแบบพลวัตจะทำงานได้ดี แต่ยังคงพบความคลาดเคลื่อนเฉลี่ยประมาณ 1--2 มม. ในกรณีที่ภาพถ่ายมีแสงสว่างไม่เพียงพอหรือมีการเอียงของมุมกล้องที่รุนแรง สาเหตุหลักมาจากการใช้ Bounding Box (รูปสี่เหลี่ยม) เป็นตัวแทนของขอบเขตวงกลม ซึ่งหากจานเพาะเชื้อเอียง Bounding Box จะวัดขนาดตามแกนที่กว้างที่สุดของวงรี ทำให้ค่าที่ได้เบี่ยงเบนไปจากเส้นผ่านศูนย์กลางจริงเล็กน้อย

\subsection{ความหลากหลายของสายพันธุ์เชื้อ (Dataset Diversity)}

ชุดข้อมูลจำนวน 1,884 ภาพที่ใช้ฝึกสอนส่วนใหญ่เป็นเชื้อ \textit{S. aureus} และ \textit{E. coli} แม้โมเดลจะมีความแม่นยำสูงมากกับเชื้อกลุ่มนี้ แต่เชื้อแบคทีเรียบางชนิดที่มีรูปแบบการเจริญเติบโตเฉพาะตัว เช่น \textit{Proteus} ที่มีการเจริญแบบ Swarming หรือ \textit{Klebsiella} ที่มีลักษณะเหนียวข้น (Mucoid) อาจสร้างขอบเขตโซนที่แปลกไปจากที่โมเดลเคยเรียนรู้ ซึ่งอาจส่งผลต่อความแม่นยำในการตีกรอบของระบบ

\subsection{ขนาดกลุ่มตัวอย่างในการประเมินโดยผู้ใช้ที่จำกัด (Small User Evaluation Sample)}

การประเมินความพึงพอใจของผู้ใช้งานในโครงงานนี้ดำเนินการกับผู้เข้าร่วมเพียง 5 ท่าน ซึ่งถือว่าเป็นขนาดกลุ่มตัวอย่างที่เล็กมากทางสถิติ ผลการประเมินเชิงคุณภาพและเชิงปริมาณที่ได้ จึงยังไม่สามารถใช้สรุปอ้างอิงเพื่อเป็นตัวแทนประชากรผู้ใช้งานระบบในวงกว้างได้อย่างมีนัยสำคัญ การศึกษาความพึงพอใจในอนาคตควรขยายกลุ่มตัวอย่างให้ครอบคลุมนักเทคนิคการแพทย์จากหลายสถานพยาบาลและหลายภูมิภาค

\subsection{ข้อจำกัดของ Hough Transform Baseline เนื่องจาก Data Leakage}

ค่าพารามิเตอร์ที่เหมาะสมที่สุดของวิธี Hough Circle Transform ถูกปรับจูนบนชุดข้อมูลทดสอบ (Test partition) ทั้งหมด ส่งผลให้ตัวเลขประสิทธิภาพที่รายงาน (MAE $>$ 5 มม.) เป็นกรณีที่ดีที่สุด (Best-case) สำหรับวิธีนี้ภายใต้เงื่อนไขของชุดข้อมูลนี้ ในสภาพแวดล้อมที่ไม่เคยเห็นมาก่อน ประสิทธิภาพของ Hough Transform น่าจะต่ำกว่าตัวเลขที่รายงาน ซึ่งหมายความว่าข้อได้เปรียบของโมเดล Deep Learning ที่แสดงในบทที่ 4 อาจมีนัยสำคัญมากกว่าที่ปรากฏจริง

\subsection{ขอบเขตการรองรับชนิดยาปฏิชีวนะของระบบ OCR (OCR Coverage Limited to Tested Antibiotics)}

ระบบ OCR และฐานข้อมูลจุดตัด (Breakpoint) ในปัจจุบันได้รับการทดสอบและตรวจสอบความถูกต้องเฉพาะกับยาปฏิชีวนะ 4 ชนิด ได้แก่ Penicillin (P), Oxacillin (OX), Chloramphenicol (C) และ Kanamycin (CN) ตามที่มีในชุดข้อมูลของโครงงาน ในห้องปฏิบัติการคลินิกจริงซึ่งอาจใช้ยาปฏิชีวนะมากกว่า 20--30 ชนิดพร้อมกัน ระบบจำเป็นต้องได้รับการขยายฐานข้อมูลและทดสอบเพิ่มเติมก่อนนำไปใช้จริงในขนาดเต็ม

% ─────────────────────────────────────────────────────────────
\section{งานในอนาคต (Future Work)}
\label{sec:future_work}
% ─────────────────────────────────────────────────────────────

แนวทางการพัฒนาและปรับปรุงระบบในระยะต่อไป ประกอบด้วยประเด็นสำคัญ ดังนี้:

\subsection{การศึกษาความถูกต้องทางคลินิกขนาดใหญ่ (Large-Scale Clinical Validation Study)}

ก่อนที่จะนำระบบไปใช้ในการวินิจฉัยเพื่อสั่งจ่ายยาจริง จำเป็นต้องมีการศึกษาความถูกต้องแบบปกปิดข้อมูล (Blinded Validation Study) อย่างเป็นทางการ โดยเปรียบเทียบผลการวัดขนาด ZOI ระหว่างระบบอัตโนมัติกับมติจากนักเทคนิคการแพทย์ผู้เชี่ยวชาญอย่างน้อย 3 ท่าน บนกลุ่มตัวอย่างเชื้อทางคลินิกที่หลากหลาย เพื่อประเมินความสอดคล้องเชิงประเภท (Categorical Agreement: CA) ตามมาตรฐาน ISO 20776-2 \cite{iso20776_2} และเกณฑ์ยอมรับของ CLSI MM17 \cite{clsi_mm17} ซึ่งครอบคลุมการวิเคราะห์อัตราความผิดพลาดร้ายแรง (Major Errors) และความผิดพลาดเล็กน้อย (Minor Errors) เพื่อเป็นข้อมูลประกอบการขอรับรองมาตรฐานเครื่องมือแพทย์ในลำดับต่อไป

\subsection{โมดูลการวิเคราะห์ระบาดวิทยาและแนวโน้มการดื้อยา (Epidemiology and Resistance Trend Module)}

การขยายขีดความสามารถของระบบในระยะยาวคือ การเพิ่มโมดูลวิเคราะห์ข้อมูลทางระบาดวิทยา (Epidemiology) เพื่อรวบรวมและวิเคราะห์สถิติการดื้อยาจากข้อมูลการตรวจทั้งหมดในฐานข้อมูล ระบบสามารถแสดงผลเป็นแดชบอร์ดแนวโน้มการดื้อยาตามชนิดเชื้อ ช่วงเวลา หรือแผนกผู้ป่วย ซึ่งจะเปลี่ยนระบบจากเพียงเครื่องมือวัดให้กลายเป็นระบบบริหารจัดการข้อมูลห้องปฏิบัติการที่มีคุณค่าต่อการเฝ้าระวังเชื้อดื้อยาในระดับสาธารณสุข สอดคล้องกับแนวทางการดำเนินงานขององค์การอนามัยโลก (WHO) \cite{who_amr2022}
